\begin{document}

\section{関連研究}

\subsection{検知後調査としての行動列復元}

エンドポイント調査では，Windows Securityログ，Sysmon，EDR telemetry，DNS，ブラウザ履歴など，粒度や記録目的の異なるログを横断する．
Sysmonはプロセス生成，親プロセス，コマンドライン，ネットワーク接続，レジストリ操作などを記録できる\cite{sysmon}．
EDRは検知理由やプロセス情報をまとめたalert summaryを提供する場合がある．
これらは調査に有用だが，alert summaryは観測された全行動を代表するものではない．
したがって，最終出力はアラート文面ではなく，ログ上で確認できる行動主張に限定する必要がある．

本研究の課題は，異常検知モデルの精度評価ではない．
DeepLogのような系列異常検知研究は，ログ系列から異常を検出・診断することに焦点を置く\cite{deeplog}．
一方，本研究は検知後のSOC調査において，調査者が検証可能な証跡付き行動列を復元できるかを問う．
この違いは評価単位にも表れる．
本研究では自然文の類似度ではなく，主体，操作，対象，コマンドライン，証跡行などの必須itemを採点する．

\subsection{LLMエージェントによるログ探索}

LLMエージェントは，初期手がかりから仮説を立て，外部環境への問い合わせ結果に応じて次の探索を変更できる．
ReActは推論と行動を交互に進める枠組みを示しており\cite{react}，ログ調査でも，時刻範囲，プロセス木，ファイル・レジストリ対象を段階的に確認する流れと親和性がある．

セキュリティ分野では，CLOUSEAUが階層型multi-agentによる自律的な攻撃調査を提案している\cite{clouseau}．
CLOUSEAUは，raw incident logsをLLMが扱いやすい構造化形式へ変換し，QA agentを通じてログを自然言語的に検索可能にする．
さらに，Investigator agentとChief Inspector agentが調査leadの生成，証跡収集，attack storyの統合を行う．
本研究はこの設計思想に基づくが，評価対象を「攻撃story全体の復元」から「SOCアラートやprocess/time clue周辺の良性を含む行動chainの証跡付き復元」へ移す．
そのため，最終判定の正否ではなく，観測ログで支えられた行動itemと過剰主張を採点する．

\subsection{ATLAS/ATLASv2の位置づけ}

ATLASは，audit logに基づいて攻撃ストーリを復元するため，因果依存グラフと系列学習を組み合わせた攻撃調査支援手法である\cite{atlas}．
ATLASv2は，ATLASの攻撃シナリオを再現しつつ，通常業務活動とSysmonおよびCarbon Black Cloud由来のログを拡充したデータセットとして報告されている\cite{atlasv2}．
本研究では，ATLASv2を「攻撃検知済みの答えを読むデータ」ではなく，「SOCでアラートやtelemetryを受け取った後，端末上の行動を証跡付きで説明する能力を測るデータ」として利用する．
この立ち位置により，本研究の対象インシデント分析は，侵害全体の帰属や最終判定ではなく，アラート周辺の端末行動復元であると明確化される．

\end{document}
